\documentclass[10pt,a4paper]{amsart}
\usepackage[margin=19mm]{geometry}
\usepackage{amsmath,amssymb,mathtools}
\usepackage[scr=rsfs,cal=euler]{mathalfa}
\usepackage{xfrac}
\usepackage[unicode,hidelinks,bookmarks=false]{hyperref}
\newcommand{\ie}{i.e.\ }
\newcommand{\cf}{cf.\ }
\newcommand{\resp}{respectively\ }
\newcommand{\Subsection}{\subsection}
\providecommand{\nobreakdash}{\nobreak\hbox{-}}
\setlength{\emergencystretch}{2em}
\multlinegap0pt
\newtheorem{lemma}{Lemma}
\newcommand{\R}{\mathbb{R}}
\newcommand{\pbra}[1]{{\left( #1 \right)}}
\newcommand{\vabs}[1]{{\left\| #1 \right\|}}
\title{Optimal Sparsification of Gaussian Processes}
\author{Shivam Nadimpalli}
\date{Source: arXiv:2606.19763v1 (CC BY 4.0)}
\begin{document}
\maketitle

\noindent\emph{Source and licence.} Optimal Sparsification of Gaussian Processes. Version arXiv:2606.19763v1 (CC BY 4.0), distributed by arXiv under the Creative Commons Attribution 4.0 International licence, http://creativecommons.org/licenses/by/4.0/, as recorded in the arXiv licence metadata for this exact version and shown on the abstract page https://arxiv.org/abs/2606.19763v1. Full licence notice: \texttt{licenses/arxiv-2606.19763-CC-BY-4.0.txt}.\par\smallskip

\noindent\textit{Editorial scope.} Counted unit: the article's lemma lem:l2-chain (source0.tex lines 3064--3116). The article's complete original proof of it is delivered with it (source0.tex lines 3118--3142, one \texttt{proof} environment of 25 lines). No mathematics is written, repaired or re-derived: the statement and the proof are the article's own lines, in the article's own notation, and no retained line is substituted or re-flowed.  No further source-local context is needed: every cross-reference of every retained line resolves inside the two delivered spans.  Nothing was omitted from any retained span, so the package carries no omission record.  No retained line contains a citation, so the wrapper carries no bibliography and no bibliography entry.  No retained line contains a figure environment, an image-inclusion command or drawing code, so no image is delivered and the package declares no asset dependency.  Editorial formatting only, in addition: an AMS \texttt{amsart} preamble that restates the article's own declarations and macros used by the retained lines, namely the environments \texttt{lemma} on separate counters, together with the standard packages those lines need.  The retained environments print in this document as Lemma 1 in order of appearance, whereas the article prints the same items with the numbering of its own full document.  The complete native source of the article as distributed in the arXiv e-print for this exact version is bundled unchanged as \texttt{sources/203169/source0.tex}.
\begin{lemma}[Softmax paths as $L^2$-valued absolutely continuous curves]
\label{lem:l2-chain}
	Let $T$ be finite.
	Suppose $v_t:[0,1]\to\R^n$ and $a_t:[0,1]\to\R$ are absolutely continuous and that, for constants $L,M<\infty$,
	\[
		\|v_t'(\theta)\|_2\leq  L\,,
		\qquad
		|a_t'(\theta)|\leq  M
	\]
	for every $t\in T$ and for a.e.~$\theta\in[0,1]$.
	Fix $\beta>0$, and set
	\[
		H_\theta(g)
		=
		\frac{1}{\beta}
		\log\pbra{
		\sum_{t\in T}
		\exp\pbra{\beta\pbra{g\cdot v_t(\theta)+a_t(\theta)}}
		}\,.
	\]
	Then $\theta\mapsto H_\theta$ is absolutely continuous as a map from $[0,1]$ into $L^2$.
	Moreover, for a.e.~$\theta$, the $L^2$-derivative is represented by
	\[
		\frac{d}{d\theta}H_\theta(g)
		=
		\sum_{t\in T}
		\omega_t^\theta(g)
		\pbra{
		g\cdot v_t'(\theta)
		+
		a_t'(\theta)
		}\,,
	\]
	where
	\[
		\omega_t^\theta(g)
		:=
		\frac{
		\exp\left(\beta(g\cdot v_t(\theta)+a_t(\theta))\right)
		}{
		\sum_{t'\in T}
		\exp\left(\beta(g\cdot v_{t'}(\theta)+a_{t'}(\theta))\right)
		}\,.
	\]
%	Consequently,
%	\[
%		\|H_1-H_0\|_{L^2}
%		\leq
%		\int_0^1
%		\vabs{\frac{d}{d\theta}H_\theta}_{L^2}
%		\,d\theta\,.
%	\]
\end{lemma}

\begin{proof}
	For each fixed $g$, the derivative formula follows from the ordinary chain rule for absolutely continuous functions, since log-sum-exp is smooth.
	Moreover, because the $\omega_t^\theta(g)$'s are non-negative and sum to one,
	\[
		\left|
		\frac{d}{d\theta}H_\theta(g)
		\right|
		\leq 
		\sum_{t\in T}
		\omega_t^\theta(g)
		\pbra{|g\cdot v_t'(\theta)|+|a_t'(\theta)|}
		\leq 
		L\|g\|_2+M\,.
	\]
	The bound $L\|g\|_2+M$ is in $L^2(\gamma_n)$ and is independent of $\theta$. 
	Therefore the difference quotients are dominated in $L^2$, the pointwise derivatives are Bochner integrable, and dominated convergence allows us to pass from pointwise absolute continuity to absolute continuity as an $L^2$-valued curve. 
	Thus, for every $0\leq \theta_1\leq \theta_2\leq 1$,
	\[
		H_{\theta_2}-H_{\theta_1}
		=
		\int_{\theta_1}^{\theta_2}
		\frac{d}{d\theta}H_\theta\,d\theta\,.
	\]
	This proves both absolute continuity and the claimed derivative formula. 
\end{proof}

\end{document}
